\section{Stage 233: Crossing-vs-collapse Goldilocks window}
\label{sec:app-part08:crossing-vs-collapse-goldilocks-window}
\label{stage:233}
\stagefield{Verification}{SymPy audit: \StageFile{scripts/moving_throat_pde_stage233_crossing_vs_collapse_goldilocks_window_compiler_from_the_stage231_event_chain_and_relaxed_wall_timescale_closure_sympy_audit.py}.  Mathematica audit: none yet.}


Stage 233 turns the dynamic branch into a survival inequality.  The exact path-time integral from the Stage-231 energy law is
\begin{equation}
\label{eq:app-part08-stage233-path-time}
T_{\rm traj}(E;r_a\to r_b)
=
\sqrt{\frac{m_s}{2}}
\int_{r_b}^{r_a}\frac{dr}{\sqrt{E-V(r)}}.
\end{equation}
The characteristic-width closure used by Session III compresses the active barrier region to
\begin{equation}
\label{eq:app-part08-stage233-cross-time}
t_{\rm cross}(E)
=
\lambda_{\rm eff}\sqrt{\frac{m_s}{2(E-V_{\rm peak})}}.
\end{equation}
This is a reduced timescale closure, not a replacement for the exact path integral.

The active dressing leg uses the local unstable normal form
\begin{equation}
\label{eq:app-part08-stage233-v-leg}
\mu_\eta\ddot V=g_{UV}\chi_{\rm peak}V.
\end{equation}
Hence
\begin{equation}
\label{eq:app-part08-stage233-collapse-time}
t_{\rm collapse}=
\sqrt{\frac{\mu_\eta}{g_{UV}\chi_{\rm peak}}}.
\end{equation}
The survival ratio is
\begin{equation}
\label{eq:app-part08-stage233-survival-ratio}
\mathcal S(E)=\frac{t_{\rm cross}(E)}{t_{\rm collapse}}
=
\lambda_{\rm eff}
\sqrt{\frac{m_sg_{UV}\chi_{\rm peak}}{2\mu_\eta(E-V_{\rm peak})}}.
\end{equation}
The lower safe edge is \cref{eq:app-part08-main-goldilocks-edge}, and the speed-space identity is \cref{eq:app-part08-main-goldilocks-speed}.

If \(\mu_\eta=\alpha m_s\), the edge becomes
\begin{equation}
\label{eq:app-part08-stage233-heavy-edge}
E_{\rm edge}=V_{\rm peak}+\frac{g_{UV}\chi_{\rm peak}\lambda_{\rm eff}^{2}}{2\alpha},
\end{equation}
so simultaneous scaling of particle mass and wall inertia cancels out of the edge.  The real control packet is
\begin{equation}
\label{eq:app-part08-stage233-control-packet}
(\lambda_{\rm eff},\chi_{\rm peak},g_{UV},\mu_\eta/m_s).
\end{equation}
With the trigger-width specialization,
\begin{equation}
\label{eq:app-part08-stage233-trigger-width}
\lambda_{\rm eff}=\lambda_{\rm th}(r_{\rm turn})=
\left|\frac{V(r_{\rm turn})}{V'(r_{\rm turn})}\right|.
\end{equation}
The benchmark edge and sensitivity statements are therefore direct consequences of the dependence
\begin{equation}
\label{eq:app-part08-stage233-edge-sensitivity}
E_{\rm edge}-V_{\rm peak}\propto\lambda_{\rm eff}^{2}\chi_{\rm peak}.
\end{equation}

Within this closure the safe window is one-sided because \(t_{\rm collapse}\) is energy-independent and \(t_{\rm cross}(E)\) decreases monotonically.  Any upper edge requires an additional high-energy failure mechanism outside Stage 233.
